\chapter{主观题}
\begin{enumerate}
  \item 找主体、看行为、分阶段、辨诉求、摆事实、讲规则。
  \item 题目类型。
        \begin{enumerate}
          \item 常识型。
          \item 直接引用法条型。写全。
                \begin{enumerate}
                  \item \al{起诉期 6 个月，但是另有规定的除外}。
                  \item \al{许可的延期。提前 30 天申请、作出答复、没有答复视为同意}。
                  \item \al{立案、不立案、7 天决定、补材料}。
                \end{enumerate}
          \item 评价、定性、说理题。\zx{是否成立。结论、理由、正确做法}。\bs{一定要说理}。\al{属于行政强制措施，因为 xxx}。
          \item 轻度结合题干。理解性。\al{是否可诉}。
          \item 措施类。只需要答最前头的，后续各种情况处理不需要明确。\al{材料告知补正书 \yjt \scx{符合的 \yjt 受理}}。
          \item 完美融合。
        \end{enumerate}
  \item 常考知识点。重心在诉讼法。
  \item \bs{专门的司法解释}。
        \begin{enumerate}
          \item 2 个程序。\zd{行政诉讼、行政复议}。
          \item 3 个行为。\zd{许可、行政协议、信息公开}。分程序法、实体法。
          \item 4 个制度。\zd{撤诉}、行政\zd{诉讼证据}、\zd{座谈会}（法律适用）、被告资格。
          \item 司法解释的 4 个子制度。\zd{负责人}出庭、复议机关作为\zd{共同被告}、民事争议的\zd{一并}审理、规范性文件的\zd{一并}审查。
        \end{enumerate}
\end{enumerate}
